\section{数据与评价口径}\label{sec:data}

\subsection{数据来源与筛选}

实验使用 Face\_segment 数据集及其原始多类别像素级标签。我们从 \texttt{test/real\_image/2} 目录中筛选出外观符合亚洲人脸特征的样本，并保留对应的像素标签。初选共得 163 对图像与标签；在配对、尺寸与标签取值核查中，发现样本 00583 与 00584 的原图完全一致而标签存在差异，经核对保留 00583、剔除 00584，最终确立 162 对有效样本。所有原图均为 $512\times512$ 的 RGB 彩色图像，标签为单通道类别图，包含 0--18 共 19 个语义类别。通过对原图、彩色编码标签与半透明叠加图的逐张核查，确认图像与标注配对准确、边界清晰。

\subsection{数据集划分与去重检查}

我们将 162 对样本按编号排序后，使用固定随机种子 42 按约 $6:2:2$ 的比例划分为训练集 98 张、验证集 32 张与测试集 32 张。在划分前，采用 64 位感知哈希（dHash，汉明距离阈值 $\le 8$）对全集进行近重复筛查，并对筛出的 35 对候选图像逐一复核，确认不存在同一照片裁剪或缩放重复。训练集用于统计肤色模型参数与五官几何先验，验证集用于算法消融与超参数选择，测试集用于检验基础流程的泛化表现。

\subsection{两套标签映射：建模用皮肤，评价用整脸}

本实验在“肤色统计建模”与“完整脸部区域评价”两个环节采用不同的标签子集：肤色建模仅提取皮肤和鼻子像素；而在评价完整脸部分割效果时，正类除皮肤和鼻子外，还包含双眼、双眉、口腔和上下唇。眼镜区域设为忽略类别，其余类别均归为负类。完整映射关系见表~\ref{tab:labels}。

\begin{table}[H]
\centering\small
\caption{原始标签在颜色统计与脸部评价中的映射规则}\label{tab:labels}
\begin{tabularx}{\textwidth}{l l X}
\toprule
用途 & 标签编号 & 对应区域 \\
\midrule
肤色统计 & 1、2 & 皮肤、鼻子 \\
脸部评价正类 & 1、2、4--7、10--12 & 皮肤、鼻子、双眼、双眉、口腔、上下唇 \\
评价忽略区域 & 3 & 眼镜 \\
脸部评价负类 & 其余类别 & 背景、双耳、头发、帽子、饰品、脖颈、衣服 \\
\bottomrule
\end{tabularx}
\end{table}

这一差异化设计符合人脸图像的物理特性：眼睛、眉毛和嘴部的颜色通常偏离皮肤色度分布，若将其纳入肤色统计会污染均值与方差估计；但完整的人脸掩膜应当涵盖五官区域，因此颜色阈值筛选在五官位置留下的天然空洞，恰好成为后续连通域填孔与形态学修补的作用对象。实现上通过集合包含关系生成二值真值掩膜与有效评测掩膜，见代码~\textbf{2-1}。

\noindent\textbf{代码 2-1}\quad 脸部真值与有效区域的生成（dataset.py）
\begin{lstlisting}
POSITIVE_IDS = [1, 2, 4, 5, 6, 7, 10, 11, 12]
IGNORE_IDS = [3]

def binary_labels(labels):
    face = np.isin(labels, POSITIVE_IDS).astype(np.uint8)
    valid = np.logical_not(np.isin(labels, IGNORE_IDS)).astype(np.uint8)
    return face, valid
\end{lstlisting}

代码中 \texttt{face} 标记完整脸部正类，\texttt{valid} 排除眼镜区域；评价时仅在 \texttt{valid} 为 1 的像素上统计混淆矩阵。

\begin{figure}[H]
\centering
\includegraphics[width=\textwidth]{part_a_label_review_00546.png}
\caption{样本 00546 的原图、原始多类别标签及半透明叠加图。}
\label{fig:label}
\end{figure}

图~\ref{fig:label} 展示了原始标签对皮肤、鼻子、眼睛、眉毛和嘴部的精细划分，同时可以看到头发和脖颈虽与脸部紧邻，但在评测定义中属于负类。明确这一语义边界有助于区分“面部内部五官漏检”与“面部外围头发脖颈误检”。

\subsection{评价指标体系}

在测试集 32 张图像共 8,388,608 个有效像素中，脸部正类像素为 2,150,722 个，占 25.64\%，背景与非脸部负类占 74.36\%。由于正负类比例显著失衡，若将全图所有像素均判为背景，像素准确率（Accuracy）仍高达 74.36\%，但此时 F1 为 0；反之，若将全图均判为脸部，F1 仅为 40.80\%（见表~\ref{tab:trivial}）。

\begin{table}[H]
\centering\small
\caption{平凡预测与基础方法在测试集上的指标对照，单位为百分比}\label{tab:trivial}
\begin{tabular}{lrrr}
\toprule
预测方式 & Accuracy & Precision & F1 \\
\midrule
全判为背景 & 74.36 & -- & 0 \\
全判为脸部 & 25.64 & 25.64 & 40.80 \\
纯阈值基线（$k=1.50$） & 76.53 & 52.91 & 62.66 \\
实验三基础方案 & 83.94 & 66.43 & 70.69 \\
\bottomrule
\end{tabular}
\end{table}

因此，本报告以 F1 与交并比（IoU）作为核心综合指标，辅以精确率（Precision）和召回率（Recall）刻画误检与漏检之间的权衡，Accuracy 仅作参考。多图汇总采用两种口径：一是先累加全部图像的 TP、FP、TN、FN 再计算整体指标（全局指标，Global），二是对每张图像单独计算指标后再取算术平均（逐图均值，Per-image Mean）。正文默认报告全局指标，并在消融与稳定性分析中同时列出逐图均值。各指标定义如式~\eqref{eq:metrics} 所示：
\begin{equation}
\begin{aligned}
\mathrm{Precision}&=\frac{TP}{TP+FP}, &
\mathrm{Recall}&=\frac{TP}{TP+FN},\\
\mathrm{F1}&=\frac{2TP}{2TP+FP+FN}, &
\mathrm{IoU}&=\frac{TP}{TP+FP+FN}.
\end{aligned}
\label{eq:metrics}
\end{equation}
